\documentclass[10pt]{article}
\usepackage[margin=1in]{geometry}
\usepackage{amsmath,amssymb,amsthm,booktabs,graphicx,microtype}
\usepackage[authoryear,round]{natbib}
\usepackage[hidelinks]{hyperref}
\usepackage{xcolor}
\usepackage{enumitem}
\usepackage{multirow}
\usepackage{siunitx}

\title{Post-Selection Certification for Safe Context Selection under Downstream-Model Uncertainty}
\author{Tao Ran\\Anhui University, Hefei, Anhui, China\\ORCID: 0009-0002-4449-8982\\Corresponding author: 2807215626@qq.com}
\date{}

\newcommand{\PSC}{\textsc{PSC-BS}}
\newcommand{\eo}{\mathrm{EO}^2}
\newcommand{\AUC}{\mathrm{AUC}}
\newcommand{\LL}{\mathrm{LL}}
\newtheorem{theorem}{Theorem}[section]
\newtheorem{proposition}[theorem]{Proposition}
\newtheorem{corollary}[theorem]{Corollary}
\newtheorem{lemma}[theorem]{Lemma}

\begin{document}
\maketitle

\begin{abstract}
Context selection can improve group fairness in in-context learning (ICL), but aggressive selection can create post-selection failure: a context that looks favorable under proxy learners may transfer poorly to an unavailable downstream model and cause severe utility loss.
We study safe context selection under downstream-model uncertainty and introduce \PSC{} (Post-Selection Benefit-Safety Certification), a model-agnostic deployment wrapper.
A proposal stage first commits to a candidate context and a deployable safety anchor without accessing an independent certification split.
Only then is that split revealed.
Three heterogeneous inexpensive certifiers compare candidate against anchor; the candidate is deployed only when every certifier satisfies frozen AUC and log-loss tolerances and the certifier ensemble shows average fairness benefit.
Otherwise, \PSC{} falls back to the anchor.
Independent sample splitting preserves fixed-pair deviation validity conditional on the prior search, yielding anchor-relative tolerance-plus-estimation-radius bounds for the certifier family.
We then test transfer empirically on unseen learners.
In frozen fresh confirmation across nine fairness tasks, five outer splits, four context budgets, four held-out classical learners, TabPFN-3.5, and TabICLv2, all preregistered CPU and GPU gates pass.
On the two unseen tabular foundation models, \PSC{} reduces severe AUC-loss events from 7.78\% for the uncertified selected candidate to 0.28\%, while reducing mean smooth equalized-odds disparity relative to the safety anchor and preserving mean utility.
Ablations identify the utility guard, heterogeneous certification, and fallback as the main downside-control mechanisms.
One genuine unseen-model miss remains; the method controls rather than eliminates tail risk.
\end{abstract}

\noindent\textbf{Keywords:} in-context learning; post-selection certification; fairness; downstream-model uncertainty; tabular foundation models; selective deployment

\section{Introduction}

In-context learning makes the context set part of the prediction algorithm.
For tabular foundation models, changing a few context rows can materially change both predictive utility and group fairness.
Recent work has therefore explored fairness-aware demonstration selection for language models and tabular foundation models \citep{hu2024strategic,kenfack2026fairtabicl,chhikara2025smite}.
A central difficulty remains underexplored: the learner used to score or validate a context during selection need not be the learner that ultimately consumes that context.

This mismatch creates \emph{downstream-model uncertainty}.
A selector may optimize a fairness proxy, a surrogate, or a validation learner and find a context that looks excellent on that proxy.
Yet the selected context is a post-selection object: it was chosen precisely because it looked favorable under the selection machinery.
When transferred to a different learner, especially a high-capacity tabular foundation model, the same context can exhibit a severe utility left tail.
Our development experiments expose this failure mode directly: fairness-oriented surrogate optimization can increase fairness wins while producing large AUC losses.
The problem is therefore not simply how to search for a fair context, but \emph{when a searched context is safe enough to deploy}.

We propose \PSC{}, a deployment wrapper that separates proposal from certification.
The proposal/search stage may be strong and even highly adaptive, but it must commit to one candidate $S^\star$ and one safety anchor $A$ before the certification split $W$ is revealed.
On $W$, a heterogeneous certifier family estimates candidate-minus-anchor changes in AUC, log-loss, and a smooth equalized-odds score.
The candidate is deployed only if it satisfies frozen utility tolerances for every certifier and shows average fairness benefit; otherwise deployment falls back to $A$.
The final downstream learner is not used in certification.
This distinguishes the setting from recent safe-ICL risk-control methods that calibrate an inference-time policy for the same known LLM relative to its zero-shot behavior; our object is a pre-deployment context decision whose transfer to an unavailable learner must be tested separately \citep{wynn2026corrupted}.

Our contributions are:
\begin{enumerate}[leftmargin=*]
    \item We formulate context selection under \emph{downstream-model uncertainty}, separating the upstream search problem from the downstream deployment-safety problem.
    \item We introduce independent post-selection benefit-safety certification with an anchor fallback, yielding a selective deployment rule rather than another context-scoring objective.
    \item We give a finite-sample framework showing why an independent certification split restores fixed-pair validity after arbitrary search on earlier data and how anchor-relative utility tolerances translate into population bounds up to estimation error.
    \item We pre-specify and cryptographically freeze the method before fresh CPU confirmation, then freeze the sequential GPU protocol before inspecting foundation-model outcomes. We evaluate on nine fairness tasks, four held-out classical learners, TabPFN-3.5, and TabICLv2; both stages pass all frozen gates.
    \item We audit tail risk rather than only averages: on unseen tabular foundation models, severe AUC losses relative to the anchor fall from 7.78\% for the raw selected candidate to 0.28\% after \PSC{}, with one genuine certification miss explicitly retained.
\end{enumerate}

\section{Related Work}

\paragraph{Fairness-aware context selection.}
Hu et al. study strategic demonstration selection for fairness in LLM ICL and use clustering and evolutionary search to curate demonstrations \citep{hu2024strategic}.
Fair-TabICL, now published in TMLR, evaluates fairness interventions for tabular foundation-model ICL, including group-balanced and uncertainty-based demonstration selection \citep{kenfack2026fairtabicl}.
SMITE jointly targets predictive and fairness errors through dynamic validation during in-context example selection \citep{chhikara2025smite}.
FairPFN takes a different route and trains a dedicated tabular foundation model for causal fairness \citep{robertson2025fairpfn}.
These works establish that context composition and tabular foundation models can be used as fairness interventions.
We do not claim first fairness-aware ICL selection or first fairness work for tabular foundation models.
Our question is downstream: once a fairness-oriented search procedure has selected a context, when is that post-selection context safe enough to deploy under a learner that was unavailable to the selector?

\paragraph{Demonstration validation and upstream selection.}
D.Va explicitly introduces demonstration validation to improve robustness and cross-model generalization \citep{zhang2025dva}; consequently, we do not claim demonstration validation itself as novel.
Coverage-based selection optimizes set-level coverage rather than independent example scores \citep{gupta2023coverage}, while recent approaches use gradient estimates or gradient matching to scale influence-based demonstration selection \citep{zzhang2025linear,jzhang2025gradient}.
We view these methods as possible proposal mechanisms.
\PSC{} is deliberately positioned \emph{after} proposal: it requires the upstream method to commit to one candidate and one fallback before the certification split is revealed.
Relative to D.Va, the distinguishing design is anchor-relative benefit-and-utility certification with an operational fallback and explicit auditing of unseen-model utility tails.

\paragraph{Post-selection inference and selective deployment.}
Sample splitting is a classical response to inference after model or variable selection \citep{kuchibhotla2022postselection}, and adaptive-data-analysis work formalizes why repeatedly reusing holdout data can destroy validity \citep{dwork2015reusable}.
Our independence theorem instantiates this familiar principle for context certification; the statistical fact of sample splitting is not the novelty claim.
The fallback also resembles selective prediction and reject-option methods, which trade coverage for reliability \citep{geifman2017selective}.
The object of rejection here is different: \PSC{} accepts or rejects an entire \emph{searched context} before an unseen downstream learner is run, rather than abstaining on individual test examples.

\paragraph{Risk-controlled safe ICL.}
A particularly close safety-oriented neighbor is Wynn et al., who use distribution-free risk control with dynamic early exit to limit degradation from corrupted or irrelevant context relative to a zero-shot baseline in language-model ICL \citep{wynn2026corrupted}.
Their calibration is built around a \emph{known} downstream LLM and selects an inference-time exit threshold; if no exit is sufficiently confident, the same model falls back to zero-shot prediction.
This work is important for our positioning because it shows that neither ``safe ICL,'' baseline-relative control, nor fallback is novel in isolation.
\PSC{} addresses a different deployment object: it certifies an entire \emph{searched context} before an unavailable downstream learner is run, uses heterogeneous external certifiers rather than internal exits of the target model, jointly checks fairness benefit and utility safety, and falls back to a predeclared deployable context rather than to zero-shot behavior.
Moreover, our frozen v2.7 rule is \emph{not} a Learn-then-Test risk-control procedure \citep{angelopoulos2025ltt}: its finite-sample theory is certifier-side fixed-pair concentration, while transfer to TabPFN-3.5 and TabICLv2 is empirical.
We therefore do not claim first risk-controlled or safe ICL; the contribution is post-selection context certification under downstream-model uncertainty.

\paragraph{Tabular foundation models.}
TabPFN established prior-data fitted transformers as strong tabular predictors \citep{hollmann2025tabpfn}; our fresh confirmation uses the recently released TabPFN-3.5 checkpoint \citep{jaeger2026tabpfn35}.
TabICL scales tabular ICL to larger context sets \citep{qu2025tabicl}; our second unseen foundation-model confirmation uses TabICLv2 \citep{qu2026tabiclv2}.
Neither TabPFN-3.5 nor TabICLv2 participates in \PSC{} search or certification.

\section{Problem Formulation}

Let a binary classification dataset consist of $(x_i,y_i,g_i)$, where $y_i\in\{0,1\}$ is the target and $g_i\in\{0,1\}$ is a protected-group indicator.
We split the available data into a context pool $P$, optional proposal calibration data $C$, a search/validation split $V$, an independent certification split $W$, and an untouched test split $T$.

Let $S\subseteq P$ denote a context of fixed budget $k$.
A proposal procedure $\mathcal A$ commits to one candidate
\begin{equation}
S^\star=\mathcal A(P,C,V),
\end{equation}
while an anchor procedure $\mathcal B$ commits to
\begin{equation}
A=\mathcal B(P,C,V).
\end{equation}
The downstream learner used on $T$ is not assumed known to $\mathcal A$, $\mathcal B$, or the certification procedure.

For predicted positive-class probabilities $p(x)$, we use a smooth score-based equalized-odds surrogate
\begin{equation}
\mathrm{EO}^2(p)
=
\sum_{y\in\{0,1\}}
\left[
\mathbb E[p(X)\mid Y=y,G=0]-
\mathbb E[p(X)\mid Y=y,G=1]
\right]^2.
\label{eq:smooth-eo}
\end{equation}
This is a smooth surrogate motivated by equalized-odds comparisons \citep{hardt2016equality}; it is not a claim that thresholded equalized odds holds at every operating point.
Smaller values indicate lower conditional score disparity.
For any metric $M$, candidate-minus-anchor differences are
\begin{equation}
\Delta M = M(S^\star)-M(A).
\end{equation}
Thus $\Delta \mathrm{EO}^2<0$ denotes a fairness improvement under the chosen surrogate, $\Delta\AUC>0$ denotes an AUC improvement, and $\Delta\LL<0$ denotes a log-loss improvement.

\section{\PSC{}: Post-Selection Benefit-Safety Certification}

\begin{figure}[t]
    \centering
    \includegraphics[width=\linewidth]{figures/fig_psc_bs_pipeline.png}
    \caption{\PSC{} separates proposal/search from independent post-selection certification. The candidate and safety anchor are fixed before $W$ is revealed. Rejection falls back exactly to the anchor; the unseen downstream learner is evaluated only on $T$.}
    \label{fig:pipeline}
\end{figure}

\subsection{Generic rule}

Let $g_1,\dots,g_K$ be inexpensive heterogeneous certifiers.
For each $g_j$, train one instance on $S^\star$ and one on $A$, then evaluate both on the same independent $W$.
Define
\begin{equation}
\Delta^{(j)}_{\AUC}, \qquad
\Delta^{(j)}_{\LL}, \qquad
\Delta^{(j)}_{\eo}.
\end{equation}

The frozen v2.7 utility certificate is
\begin{equation}
\min_j \Delta^{(j)}_{\AUC}\ge -0.01,
\qquad
\max_j \Delta^{(j)}_{\LL}\le 0.05,
\label{eq:utilitygate}
\end{equation}
and the benefit certificate is
\begin{equation}
\frac{1}{K}\sum_{j=1}^K \Delta^{(j)}_{\eo}<0.
\label{eq:benefitgate}
\end{equation}
Let $Z$ indicate that all conditions in Eqs.~\eqref{eq:utilitygate}--\eqref{eq:benefitgate} hold.
Deployment is
\begin{equation}
S_{\mathrm{deploy}}=
\begin{cases}
S^\star,&Z=1,\\
A,&Z=0.
\end{cases}
\label{eq:fallback}
\end{equation}

The utility rule is deliberately conservative across certifiers; the fairness rule requires average benefit across certifiers.
No downstream test model participates in this decision.

\subsection{Frozen instantiation}

The confirmatory implementation uses:
\begin{itemize}[leftmargin=*]
    \item candidate proposal: 16 prevalence-matched random contexts scored on $V$ by Logistic Regression and HistGradientBoosting; among validation-utility-feasible draws, choose the lowest mean smooth $\eo$;
    \item safety anchor: five prevalence-matched random contexts, with one anchor chosen on $V$ by a frozen Logistic/HistGradientBoosting/Random-Forest verifier;
    \item independent certifiers on $W$: Logistic Regression, HistGradientBoosting, and Random Forest;
    \item AUC tolerance $0.01$, log-loss tolerance $0.05$, fairness margin $0$.
\end{itemize}
The entire rule was frozen before the fresh v2.7 confirmation.

\section{Why Independence Matters}

Let $D_0=(P,C,V)$ and let $U$ contain all random seeds or algorithmic randomness used before the certification outcomes are inspected.
Define $\mathcal F_0=\sigma(D_0,U)$.
The selected candidate $S^\star$, anchor $A$, and the fitted certifier predictors (which train only on these contexts) are measurable with respect to information independent of $W$.
The certification sample satisfies
\begin{equation}
W\perp \mathcal F_0.
\end{equation}
This is a sample-splitting construction, not a new post-selection inference principle \citep{kuchibhotla2022postselection}.

\paragraph{Theorem 1 (independent post-selection pair validity).}
Suppose that for every fixed pair $(S,A)$, an estimator $\widehat R_W(S,A)$ obeys
\begin{equation}
\Pr\!\left(
\left|\widehat R_W(S,A)-R(S,A)\right|>\varepsilon_R(\delta)
\right)\le \delta.
\label{eq:fixedpair}
\end{equation}
If $(S^\star,A)$ is fixed before $W$ is inspected, then
\begin{equation}
\Pr\!\left(
\left|\widehat R_W(S^\star,A)-R(S^\star,A)\right|>\varepsilon_R(\delta)
\mid \mathcal F_0
\right)\le \delta
\end{equation}
almost surely, and therefore the same bound holds unconditionally.

\paragraph{Proof.}
Condition on any realization of $\mathcal F_0$.
The random selected pair becomes a fixed pair $(s,a)$.
By independence, conditioning on $\mathcal F_0$ does not change the law of $W$.
Equation~\eqref{eq:fixedpair} therefore applies to $(s,a)$ with conditional failure probability at most $\delta$.
Taking expectation over $\mathcal F_0$ gives the unconditional statement. $\square$

\paragraph{Why search on $W$ is different.}
If $M$ candidate contexts are scored and one is selected using the same $W$, validity generally requires simultaneous rather than fixed-pair control.
For a bounded empirical-mean statistic, a union bound with Hoeffding's inequality gives
\begin{equation}
\Pr\!\left(
\sup_{S\in\mathcal C}|\widehat R_W(S)-R(S)|>\epsilon
\right)
\le 2M e^{-2n_W\epsilon^2},
\end{equation}
so a representative radius scales as
\begin{equation}
\epsilon=\sqrt{\frac{\log(2M/\delta)}{2n_W}}.
\end{equation}
\PSC{} does not pay this \emph{search-over-$W$} factor for the committed pair because proposal occurs on earlier data.
This argument does not prevent adaptive analysis of other holdouts in general; it only describes the frozen one-shot certification protocol used here \citep{dwork2015reusable}.

\subsection{Anchor-relative utility bounds}

Let $R$ be a risk where smaller is better and define $\Delta_R=R(S^\star)-R(A)$.
If
\begin{equation}
|\widehat\Delta_R-\Delta_R|\le \epsilon_R
\end{equation}
with probability at least $1-\delta$, then acceptance under $\widehat\Delta_R\le \tau$ implies
\begin{equation}
\Delta_R\le \tau+\epsilon_R.
\label{eq:riskcor}
\end{equation}
For AUC, acceptance under $\widehat\Delta_{\AUC}\ge-\tau_{\AUC}$ analogously implies
\begin{equation}
\Delta_{\AUC}\ge-\tau_{\AUC}-\epsilon_{\AUC}.
\label{eq:auccor}
\end{equation}
The frozen v2.7 gate uses point tolerances without an extra statistical margin.
Equations~\eqref{eq:riskcor}--\eqref{eq:auccor} therefore justify a \emph{tolerance plus estimation radius} interpretation; they do not retroactively turn the empirical 0.01/0.05 tolerances into exact population guarantees.

\subsection{Metric-specific deviation control}

The three certified metrics require different arguments.
For log-loss, predictions are clipped to $[\eta,1-\eta]$ with $\eta=10^{-8}$, so per-example log-loss lies in $[0,B]$ where $B=-\log\eta$.
The paired candidate-minus-anchor log-loss contribution therefore lies in $[-B,B]$, and Hoeffding gives the valid fixed-pair radius
\begin{equation}
\epsilon_{\LL}(\delta)=B\sqrt{\frac{2\log(2/\delta)}{n_W}}.
\end{equation}

Empirical AUC is a bounded two-sample U-statistic rather than an iid empirical mean.
Accordingly, we rely on standard U-statistic concentration for each fixed predictor and combine candidate/anchor deviations by a union bound \citep{clemencon2008ranking}.
We leave the radius as $\epsilon_{\AUC}(\delta;n_+,n_-)$ to avoid pretending that class imbalance has the same concentration behavior as a simple sample mean.

For smooth $\eo$, let
\begin{equation}
d_y(f)=\mathbb E[f(X)\mid Y=y,G=0]-\mathbb E[f(X)\mid Y=y,G=1].
\end{equation}
Conditional on the four realized cell counts $n_{yg}$, Hoeffding controls each bounded score mean.
If all eight candidate/anchor group-conditional means are simultaneously within radii $r_{yg}$ of their expectations, then $|a^2-b^2|\le2|a-b|$ for $a,b\in[-1,1]$ yields
\begin{equation}
|\widehat\Delta_{\eo}-\Delta_{\eo}|
\le 4\sum_{y,g} r_{yg}.
\label{eq:eobound}
\end{equation}
One valid choice is
\begin{equation}
r_{yg}=\sqrt{\frac{\log(16/\delta)}{2n_{yg}}}.
\end{equation}
The bound is intentionally conservative; its role is to make the finite-sample logic explicit, not to claim that the deployed point gate is a tight confidence procedure.

For $K$ certifiers and multiple required metric events, simultaneous validity follows by allocating the desired failure probability across the relevant events with a union bound.
The utility gate uses a worst-certifier intersection, whereas the fairness gate targets the average certifier $\Delta\eo$.

\subsection{Fallback and downstream-model uncertainty}

Let $L_h$ denote a bad deployment event for an unseen learner $h$ relative to the anchor, such as $\Delta_h\AUC<-0.05$.
Because rejection deploys the anchor exactly,
\begin{equation}
\Pr(L_h\text{ under }\PSC{})
=
\Pr(Z=1)\Pr(L_h\mid Z=1).
\label{eq:fallbackdecomp}
\end{equation}
This identity explains why coverage matters but does not prove that the certificate chooses the right contexts.
We therefore test that mechanism empirically, including a matched-acceptance random-gate diagnostic.

Most importantly, Theorem~1 and Eqs.~\eqref{eq:riskcor}--\eqref{eq:eobound} apply to the \emph{certifier family}.
They do not prove that an unseen learner $h$ must behave identically.
A transfer theorem would require an additional discrepancy assumption, for example
\begin{equation}
\Delta_h^U(S,A)
\le
\max_j \Delta_{g_j}^U(S,A)+\rho_h^U.
\end{equation}
We therefore separate claims cleanly: theory establishes independent post-selection validity for certified quantities; fresh held-out experiments test whether those certificates transfer to unseen learners.

\section{Experimental Protocol}

\subsection{Data}

We use nine binary fairness tasks:
seven 2023 ACS/Folktables tasks---Income, Employment, Public Coverage, Mobility, Travel Time, Health Insurance, and Income-to-Poverty Ratio---aggregated over Colorado, Michigan, Minnesota, New Jersey, and Oregon \citep{ding2021retiring};
the LSAC longitudinal bar-passage data \citep{wightman1998lsac};
and the UCI Diabetes 130-US Hospitals dataset \citep{clore2014diabetes,strack2014hba1c}.

For the ACS tasks, the protected comparison is RAC1P restricted to White-alone and Black-alone records, with race excluded from features.
For Law School, the target is first-try bar passage and the protected attribute is sex.
For Diabetes, the target is readmission within 30 days and the protected comparison is Caucasian versus African-American race.
Protected attributes are excluded from model inputs.

\subsection{Splits and preprocessing}

Each dataset uses five deterministic outer splits, stratified jointly by target and protected group.
The split sizes are $|P|=3000$, $|C|=500$, $|V|=500$, $|W|=1000$, and $|T|=2000$ for every fresh task; all filtered datasets exceed the 7000-row cap.
We evaluate context budgets $k\in\{32,64,128,256\}$.
Numeric features are median-imputed and standardized; categorical features are one-hot encoded.
The preprocessing transform is fit on $P$ only.

\subsection{Downstream models and baselines}

Fresh CPU confirmation uses four held-out model families not used in certification: ExtraTrees, XGBoost, LightGBM, and an MLP.
Fresh GPU confirmation uses TabPFN-3.5 and TabICLv2.
The eight deployed selectors evaluated on GPU are \PSC{}, the raw 16-draw selected candidate, the safety anchor, and five prevalence-random references.

The primary comparison is \PSC{} versus the safety anchor.
We additionally compare against the mean of the five prevalence-random contexts.
A \emph{severe utility tail} is predefined as
\begin{equation}
\Delta\AUC < -0.05
\end{equation}
relative to the safety anchor.
A random-baseline \emph{guarded win} requires improved smooth $\eo$, $\Delta\AUC\ge -0.01$, and $\Delta\LL\le 0.05$.

\subsection{Freshness and statistics}

The v2.7 method and CPU gates were pre-specified and frozen on 2026-09-27 before the fresh v2.7 CPU outcomes were inspected.
The frozen implementation SHA256 is
\begin{center}
\footnotesize\texttt{f02ccffec4bd20225e261bd4a2b88c803ebfbf700e8f53d6b664e42f82dc307d}
\end{center}
Development used ACS-2022 AZ/GA/MA/NC/WA; fresh CPU uses ACS-2023 CO/MI/MN/NJ/OR together with Law School and Diabetes. The GPU protocol was frozen only after CPU confirmation passed, but before TabPFN-3.5 or TabICLv2 outcomes were inspected, and it replays exactly the CPU-selected contexts. Thus the study is a sequential confirmation rather than two statistically independent preregistrations. Neither foundation model participates in search or certification.

We report paired effects, dataset/model/budget stratification, severe-tail rates, and bootstrap intervals.
Dataset-cluster bootstrap resamples entire datasets to expose cross-dataset uncertainty.


\paragraph{AI-assisted research workflow.}
OpenAI ChatGPT was used during the research workflow for code drafting and debugging, experiment-orchestration assistance, result-audit support, literature triage, and manuscript drafting/editing.
All experimental protocols, thresholds, method freezes, executed computations, source verification, interpretation, and final manuscript content were reviewed and approved by the human author(s), who take full responsibility for the work.
No AI system is listed as an author.


\section{Results}

\subsection{Fresh confirmation}

\begin{table}[t]
\centering
\caption{Fresh confirmatory results. Differences are \PSC{} minus the safety anchor unless noted. Lower $\Delta\eo$ and $\Delta\LL$ are better; higher $\Delta\AUC$ is better.}
\label{tab:fresh}
\begin{tabular}{lrr}
\toprule
 & Fresh CPU & Fresh GPU \\
\midrule
Tasks & 5760/5760 & 2880/2880 \\
Errors & 0 & 0 \\
Downstream models & 4 & 2 \\
Paired model cells & 720 & 360 \\
Acceptance & 16.11\% & 16.11\% \\
Mean $\Delta\eo$ & -0.001200 & -0.000513 \\
Mean $\Delta\AUC$ & +0.000780 & +0.001208 \\
Mean $\Delta\LL$ & -0.003554 & -0.002410 \\
Severe AUC tail & 0.417\% & 0.278\% \\
Random5 fairness win & 64.44\% & 65.00\% \\
Random5 guarded win & 52.36\% & 51.67\% \\
Random5 mean $\Delta\AUC$ & +0.011063 & +0.010031 \\
\bottomrule
\end{tabular}
\end{table}

All eight frozen CPU gates and all nine sequentially frozen GPU gates pass.
On GPU, the primary fairness effect versus anchor is negative with a paired-cell 95\% bootstrap interval $[-8.37,-2.31]\times10^{-4}$ and a dataset-task-cluster interval $[-8.77,-2.02]\times10^{-4}$.
The mean AUC change is slightly positive; its dataset-task-cluster interval crosses zero, so we interpret utility as approximately preserved rather than claiming a universal AUC gain.
The random5 fairness/guarded gates pass on their pre-specified point estimates, but their dataset-task-cluster intervals cross the 60\%/50\% thresholds; we therefore treat random5 win rates as secondary rather than the strongest robustness claim.

The scale of the fairness effect is easier to interpret in absolute terms.
Across the GPU lattice, the safety anchor has mean smooth $\eo=0.005171$ and \PSC{} has $0.004658$, a 9.9\% reduction in the aggregate mean.
Among the 58 accepted TFM model-cells, anchor smooth $\eo$ averages $0.006288$ versus $0.003105$ after deploying the certified candidate, a 50.6\% descriptive reduction; accepted cells also have mean $\Delta\AUC=+0.00750$ and $\Delta\LL=-0.01496$ on the untouched TFM test outcomes.
These accepted-subset numbers are descriptive effects, not additional selection rules.

\subsection{Tail-risk control on unseen tabular foundation models}

The raw selected candidate has 28 severe AUC-tail cells out of 360 (7.778\%).
After \PSC{} certification and fallback, only one severe-tail cell remains (0.278\%).
The observed absolute risk reduction is 7.50 percentage points; the paired-cell bootstrap 95\% interval is [5.00, 10.28] points and the dataset-task-cluster interval is [5.28, 9.44] points.

\begin{table}[t]
\centering
\caption{Fresh GPU effects by unseen downstream model. Candidate tail is the severe-tail rate before \PSC{} fallback.}
\label{tab:tfm}
\small
\begin{tabular}{lrrrrr}
\toprule
Model & $\Delta\eo$ & $\Delta\AUC$ & $\Delta\LL$ & \PSC{} tail & Candidate tail \\
\midrule
TabPFN-3.5 & -0.000363 & +0.001482 & -0.000456 & 0.000\% & 6.111\% \\
TabICLv2 & -0.000663 & +0.000934 & -0.004363 & 0.556\% & 9.444\% \\
\bottomrule
\end{tabular}
\end{table}

Both unseen models satisfy the frozen transfer gate: negative mean $\Delta\eo$ and mean $\Delta\AUC\ge-0.01$ relative to anchor.
The single remaining miss occurs on Law School, outer split 4, context budget 32, under TabICLv2: $\Delta\AUC=-0.0661$, $\Delta\eo=+0.00945$, and $\Delta\LL=+0.0819$.
This is a genuine accepted-candidate miss, not an implementation error.
Accordingly, our claim is that \PSC{} substantially \emph{reduces} or \emph{controls} severe downside incidence; it does not eliminate it.

\subsection{Mechanism ablations}

\begin{table*}[t]
\centering
\caption{GPU ablations reconstructed from frozen certification metadata and already-computed fresh GPU candidate/anchor outcomes. These analyses are post-confirmatory mechanism diagnostics.}
\label{tab:ablation}
\small
\begin{tabular}{lrrrrrr}
\toprule
Policy & Accept & $\Delta\eo$ & $\Delta\AUC$ & $\Delta\LL$ & Severe tail & Guarded win \\
\midrule
Full \PSC{} & 16.1\% & -0.000513 & +0.001208 & -0.002410 & 0.278\% & 51.67\% \\
Uncertified candidate / no fallback & 100.0\% & -0.002451 & -0.004297 & -0.008518 & 7.778\% & 55.00\% \\
Utility-only certificate & 31.7\% & -0.000426 & +0.003150 & -0.005188 & 0.278\% & 51.67\% \\
Benefit-only certificate & 55.6\% & -0.002508 & -0.003313 & -0.005157 & 5.000\% & 57.22\% \\
Single Logistic certifier & 31.7\% & -0.000765 & +0.001794 & -0.004286 & 1.389\% & 53.61\% \\
Single HistGB certifier & 31.1\% & -0.001525 & +0.002203 & -0.004777 & 0.833\% & 55.56\% \\
Single RandomForest certifier & 34.4\% & -0.001056 & +0.003062 & -0.004737 & 1.111\% & 53.61\% \\
\bottomrule
\end{tabular}
\end{table*}

Removing certification/fallback restores the 7.778\% candidate tail.
Keeping only the fairness-benefit gate leaves a 5.0\% tail, identifying the utility guard as the main safety component.
Single-certifier variants leak more tail risk than the heterogeneous intersection.
The utility-only rule matches the full method's observed tail rate but yields less aggregate fairness improvement, consistent with the fairness gate acting as benefit targeting rather than primary utility protection.

A low acceptance rate mechanically reduces failures because rejected cells revert to the anchor.
To test whether \emph{which} cells are accepted matters, we add an explicitly post-confirmatory matched-acceptance diagnostic.
Across 100,000 random gates that each accept exactly 29 of the 180 context cells, the mean severe-tail rate is 1.254\%, compared with 0.278\% for \PSC{}; mean $\Delta\AUC$ is $-0.000694$ for the matched-random gate versus $+0.001208$ for \PSC{}.
Only 8.6\% of matched-random draws have a tail count no larger than the observed \PSC{} count.
Because this diagnostic was added after confirmation, we use it as mechanistic evidence, not as a confirmatory hypothesis test.

\begin{figure}[t]
    \centering
    \includegraphics[width=\linewidth]{figures/fig_ablation_tail_rate.png}
    \caption{Severe AUC-tail rate for post-confirmatory GPU ablations.}
    \label{fig:ablation}
\end{figure}

\subsection{Sensitivity}

Reducing the certification sample from 1000 to 250 examples preserves the qualitative behavior; at 125 examples acceptance becomes more conservative and fairness benefit weakens.
Observed severe-tail incidence remains between 0 and 0.556\% over $n_W\in\{125,250,500,1000\}$.
We retain the frozen $n_W=1000$ and do not retune to the post-hoc sensitivity result.

\begin{figure}[t]
    \centering
    \includegraphics[width=\linewidth]{figures/fig_cert_size_tail.png}
    \caption{Certification-sample-size sensitivity. The zero-tail point at $n_W=500$ is a finite-sample observation, not a basis for retuning.}
    \label{fig:certsize}
\end{figure}

Candidate-count sensitivity regenerates deterministic prefixes of the same 16-draw proposal stream at $M\in\{1,2,4,8,16\}$ and reuses the unchanged certificate.
The selected candidate changes materially with $M$ (only 6.67\% of $M=1$ selections equal the $M=16$ choice), yet severe-tail incidence stays below 0.6\% for every tested $M$ on the four held-out classical model families.
There is no monotone performance trend with $M$; proposal count is a search-capacity parameter, whereas certification/fallback controls deployment risk.
The frozen $M=16$ setting is retained.

\begin{figure}[t]
    \centering
    \includegraphics[width=\linewidth]{figures/fig_candidate_count_tail.png}
    \caption{Candidate-count sensitivity on held-out classical models.}
    \label{fig:candidates}
\end{figure}

\section{Discussion}

The experiments support a distinction between \emph{search quality} and \emph{deployment safety}.
The raw fairness-oriented search finds contexts with stronger average fairness than the selective deployment rule, but it also exhibits a substantially larger unseen-model utility tail.
\PSC{} is intentionally selective: 29/180 fresh context cells deploy the searched candidate and the remainder deploy the predeclared anchor.
This behavior is analogous to a reject option at the context level rather than at the prediction level \citep{geifman2017selective}.
The pre-specified CPU acceptance gate of 10--60\% rules out the trivial always-fallback policy.

The matched-acceptance diagnostic suggests that selectivity alone does not explain the full result.
Randomly accepting the same number of contexts produces an expected severe-tail rate roughly 4.5 times larger than \PSC{} and a negative mean AUC change.
Nevertheless, because this analysis is post-confirmatory and the randomization tail probability is 0.086, we treat it as mechanism evidence rather than a new statistical claim.

The heterogeneous certifier ensemble also has a concrete role.
A single cheap model can miss failures that another model exposes; intersecting utility conditions across Logistic Regression, HistGradientBoosting, and Random Forest lowers the observed unseen-model tail relative to each single-certifier variant.
This is empirical evidence for model diversity, not a theorem that the three certifiers span every future learner.

Finally, the anchor is central and deliberately modest.
Absolute thresholds would mix task difficulty with deployment policy.
Anchor-relative certification instead asks whether the searched context is sufficiently safe and beneficial compared with a context that can actually be deployed when certification fails.
The theory is therefore only as meaningful as the chosen anchor: we do not claim that the anchor is globally optimal or absolutely fair.

\paragraph{Novelty boundary.}
We do not claim first sample splitting, first post-selection validation, first demonstration validation, first reject option, first fairness-aware context selection, or first fairness-aware tabular foundation model.
The contribution is their problem-specific combination and evidence: an unknown-downstream context-selection formulation, a committed candidate/anchor pair, independent benefit-and-utility certification, exact fallback, and explicit fresh measurement of unseen-model left-tail utility risk.

\section{Limitations}

Our evidence has several deliberate boundaries.
First, theory establishes certifier-side post-selection validity, while transfer to unseen downstream learners remains empirical unless a learner-discrepancy assumption is added.
Second, the frozen v2.7 gate uses point tolerances chosen during development (AUC 0.01, log-loss 0.05) without an additional confidence margin; they are operational thresholds rather than statistically optimal constants.
Third, the method accepts only 16.11\% of fresh context cells, so it trades deployment coverage for reliability; other applications may require a different operating point, which should be chosen before confirmation rather than tuned on these outcomes.
Fourth, the study considers binary protected-group comparisons and a smooth score-based $\eo$ surrogate; intersectional, multi-group, and threshold-specific fairness constraints remain open.
Fifth, the data are U.S.-centric. Moreover, seven of the nine task labels share the same ACS-2023 source and five states; task-level dataset-cluster bootstrap therefore should not be interpreted as nine independent data-generating domains.
Sixth, one genuine TabICLv2 certification miss remains.
Seventh, candidate-count sensitivity on alternative candidates is evaluated with classical held-out learners rather than rerunning post-hoc foundation-model inference for every $M$.
Eighth, the main fresh confirmation freezes one strong proposal selector; broader selector-swap confirmation would test wrapper generality but is not required for the frozen claim reported here.

\section{Conclusion}

We introduced \PSC{}, an independent post-selection benefit-safety certificate with a safety-anchor fallback for context selection under downstream-model uncertainty.
The statistical ingredient of sample splitting is classical; the method contribution is an operational deployment wrapper that separates fairness-oriented context search from a later anchor-relative safety decision when the eventual learner is unavailable.
Across frozen fresh CPU confirmation and a sequentially frozen TabPFN-3.5/TabICLv2 confirmation, the method passes all pre-specified gates.
The strongest empirical result is tail-focused: severe AUC-loss incidence on the unseen foundation models falls from 7.78\% for the raw selected candidate to 0.28\% after certification and fallback.
Ablations, matched-acceptance analysis, and sensitivity studies support the intended mechanism without changing the frozen v2.7 rule.
The remaining miss and the theory's transfer boundary keep the scope explicit: \PSC{} controls downside risk; it does not promise universal safety or global optimality.

\appendix
\section{Closest-Neighbor Positioning}

\begin{table*}[h]
\centering
\caption{Positioning against the closest risk-control neighbor. The distinction is not that PSC-BS uniquely uses a baseline or fallback; it is the post-selection, unknown-downstream setting and the jointly certified fairness/utility decision.}
\label{tab:closest}
\small
\begin{tabular}{p{3.1cm}p{5.3cm}p{6.1cm}}
\toprule
 & Wynn et al. (ICML 2026) & PSC-BS (ours) \\
\midrule
Primary failure mode & Harmful/irrelevant user context degrades a known LLM & Fairness-oriented context search overfits proxies and transfers poorly to an unavailable learner \\
Decision object & Inference-time early-exit threshold & Entire searched context $S^\star$ versus anchor $A$ \\
Baseline/fallback & Same LLM in zero-shot mode & Predeclared deployable anchor context \\
Calibration/certification & Distribution-free risk control / Learn-then-Test-style calibration & Independent post-selection fixed-pair certification with operational tolerances \\
Models used for the decision & The known target LLM's internal exits & Logistic, HistGB, and Random Forest; target TFM excluded \\
Quantities controlled & Loss relative to zero-shot & Worst-certifier AUC/log-loss safety plus mean smooth-EO$^2$ benefit \\
Unknown downstream model & Not the central setting & Central setting; transfer is evaluated empirically \\
Intervention after rejection & Ignore context / zero-shot prediction & Deploy anchor context \\
Theory claim & Risk control for the calibrated known-model policy & Certifier-side post-selection validity; no theorem of unseen-model transfer \\
\bottomrule
\end{tabular}
\end{table*}


\section{Formal Post-Selection Guarantees}
\label{app:formal}

This appendix states the statistical claims used by \PSC{} in a form that separates what follows from sample splitting from what remains an empirical transfer question. None of the results below upgrades the frozen v2.7 rule into a distribution-free risk-control procedure: the deployed rule uses fixed operational tolerances and no additional confidence margin.

\subsection{Setup}
Let $D_0=(P,C,V)$ collect all data and randomness used before certification, and let $\mathcal F_0=\sigma(D_0)$. The proposal and anchor are measurable with respect to $\mathcal F_0$:
\[
S^\star=\mathcal A(D_0),\qquad A=\mathcal B(D_0).
\]
The certification sample $W=(Z_1,\ldots,Z_n)$ is independent of $\mathcal F_0$. For a fixed certifier $g$ and a fixed pair $(S,A)$, let $\Delta_q(g;S,A)$ denote a population candidate-minus-anchor contrast for quantity $q$, and let $\widehat\Delta_q(g;S,A;W)$ be its estimator on $W$.

\begin{theorem}[Conditional fixed-pair validity]
\label{thm:fixedpair}
Assume that for every deterministic pair $(s,a)$ and every $\delta\in(0,1)$,
\[
\Pr\!\left(
\left|\widehat\Delta_q(g;s,a;W)-\Delta_q(g;s,a)\right|>
\varepsilon_q(\delta)
\right)\le \delta.
\]
Then the adaptively selected pair $(S^\star,A)$ satisfies
\[
\Pr\!\left(
\left|\widehat\Delta_q(g;S^\star,A;W)-\Delta_q(g;S^\star,A)\right|>
\varepsilon_q(\delta)
\mid \mathcal F_0
\right)\le \delta
\]
almost surely, and therefore the same inequality holds unconditionally.
\end{theorem}

\begin{proof}
Condition on $\mathcal F_0$. On every realized pre-certification history, $(S^\star,A)$ becomes one deterministic pair $(s,a)$. Independence implies that the conditional law of $W$ given $\mathcal F_0$ is the same law used by the fixed-pair deviation assumption. The assumed bound therefore applies conditionally. Taking expectations of the conditional failure probability over $\mathcal F_0$ gives the unconditional statement.
\end{proof}

\begin{corollary}[No search-size penalty from earlier data]
\label{cor:searchsize}
The deviation radius in Theorem~\ref{thm:fixedpair} need not contain the number $M$ of candidates searched on $V$, even if $M$ is large or the search on $V$ is adaptive, provided exactly one candidate and one anchor are committed before $W$ is inspected.
\end{corollary}

\begin{proof}
Theorem~\ref{thm:fixedpair} conditions on the complete search history. After conditioning, only one fixed pair is evaluated on $W$. A union bound over the candidate family would be required only if $W$ itself were reused to compare or select among multiple candidates.
\end{proof}

\subsection{Anchor-relative one-sided guarantees}

\begin{proposition}[Operational tolerance plus estimation radius]
\label{prop:tolerance}
Suppose an event $E$ satisfies
\[
E=\left\{ |\widehat\Delta_q-\Delta_q|\le \varepsilon_q\right\}.
\]
For a loss-like metric, acceptance under $\widehat\Delta_q\le \tau_q$ implies on $E$ that
\[
\Delta_q\le \tau_q+\varepsilon_q.
\]
For a utility-like metric such as AUC, acceptance under $\widehat\Delta_q\ge -\tau_q$ implies on $E$ that
\[
\Delta_q\ge -\tau_q-\varepsilon_q.
\]
\end{proposition}

\begin{proof}
Both statements follow directly from the triangle inequality. For the loss case, $\Delta_q\le \widehat\Delta_q+\varepsilon_q\le \tau_q+\varepsilon_q$. The AUC statement follows from $\Delta_q\ge \widehat\Delta_q-\varepsilon_q$.
\end{proof}

\begin{corollary}[Simultaneous utility bounds for heterogeneous certifiers]
\label{cor:multi}
Let $K$ certifiers be fixed before observing $W$. If AUC and log-loss deviation events for each certifier are each assigned failure probability at most $\delta/(2K)$, then with probability at least $1-\delta$, every certifier that passes the empirical frozen utility gates also satisfies
\[
\Delta^{(j)}_{\AUC}\ge -\tau_{\AUC}-\varepsilon_{\AUC,j},\qquad
\Delta^{(j)}_{\LL}\le \tau_{\LL}+\varepsilon_{\LL,j}
\]
simultaneously for all $j=1,\ldots,K$.
\end{corollary}

\begin{proof}
Apply Proposition~\ref{prop:tolerance} on the intersection of the $2K$ deviation events and use a union bound.
\end{proof}

\subsection{Metric-specific finite-sample radii}
The following radii are intentionally conservative and are provided to make the assumptions behind the generic theory explicit. They are not inserted into the frozen v2.7 deployment rule.

\begin{proposition}[Clipped log-loss]
\label{prop:logloss}
Assume class probabilities are clipped to $[\eta,1-\eta]$ with $0<\eta<1/2$, and define $B=\log(1/\eta)$. Conditional on $\mathcal F_0$, let $D_i$ be the candidate-minus-anchor per-example log-loss difference on certification point $Z_i$. Then $D_i\in[-B,B]$. If the certification observations are independent,
\[
\Pr\left(
|\widehat\Delta_{\LL}-\Delta_{\LL}|>\epsilon
\mid\mathcal F_0
\right)
\le 2\exp\left(-\frac{n\epsilon^2}{2B^2}\right).
\]
Thus one valid two-sided radius is
\[
\varepsilon_{\LL}(\delta)
= B\sqrt{\frac{2\log(2/\delta)}{n}}.
\]
\end{proposition}

\begin{proof}
Clipping bounds each binary log loss between $0$ and $B$, hence the paired difference lies in an interval of width $2B$. Hoeffding's inequality applied to the empirical mean of $D_i$ yields the stated bound.
\end{proof}

\begin{proposition}[AUC via bounded differences]
\label{prop:auc}
Condition on the certification class counts $n_1$ and $n_0$ and on $\mathcal F_0$. Assume $n_1,n_0>0$. For a fixed candidate/anchor pair, empirical AUC is a two-sample U-statistic with a kernel in $[0,1]$. For the candidate-minus-anchor AUC difference,
\[
\Pr\left(
|\widehat\Delta_{\AUC}-\Delta_{\AUC}|>\epsilon
\mid \mathcal F_0,n_1,n_0
\right)
\le
2\exp\left(
-\frac{\epsilon^2}{2(1/n_1+1/n_0)}
\right).
\]
A conservative radius is therefore
\[
\varepsilon_{\AUC}(\delta)
=
\sqrt{2\left(\frac1{n_1}+\frac1{n_0}\right)\log\frac{2}{\delta}}.
\]
\end{proposition}

\begin{proof}
Write the AUC-difference estimator as the average over all positive-negative pairs of the difference between two pairwise ranking kernels. Each pairwise difference lies in $[-1,1]$. Replacing one positive observation can change the average by at most $2/n_1$, and replacing one negative observation can change it by at most $2/n_0$. McDiarmid's bounded-differences inequality therefore gives
\[
2\exp\left(
-\frac{2\epsilon^2}{n_1(2/n_1)^2+n_0(2/n_0)^2}
\right),
\]
which simplifies to the displayed inequality. Under iid sampling within classes, the two-sample U-statistic is unbiased for the corresponding population AUC contrast. Related ranking/U-statistic concentration results are standard \citep{clemencon2008ranking}.
\end{proof}

\begin{proposition}[A conservative smooth-$\eo$ radius]
\label{prop:eo}
For context $c\in\{S^\star,A\}$, define
\[
\mu^c_{yg}=\mathbb E[p_c(X)\mid Y=y,G=g],
\qquad y,g\in\{0,1\},
\]
where $p_c(X)\in[0,1]$. Let $n_{yg}$ be the certification count in each target-group cell and $n_{\min}=\min_{y,g}n_{yg}>0$. If all eight candidate/anchor conditional means are estimated from $W$, then with probability at least $1-\delta$,
\[
|\widehat\Delta_{\eo}-\Delta_{\eo}|
\le
16\sqrt{\frac{\log(16/\delta)}{2n_{\min}}}.
\]
The right-hand side may be truncated at the trivial range bound $2$.
\end{proposition}

\begin{proof}
By Hoeffding and a union bound over the eight conditional means, with probability at least $1-\delta$ every mean has error at most
\[
\eta_0=\sqrt{\frac{\log(16/\delta)}{2n_{\min}}}.
\]
For a fixed context and target value $y$, the group-gap estimate differs from its population value by at most $2\eta_0$. Both gaps lie in $[-1,1]$, so $|\hat d_y^2-d_y^2|\le 2|\hat d_y-d_y|\le4\eta_0$. Summing over the two target values gives an $8\eta_0$ bound for one context's smooth-$\eo$ score. Taking a candidate-minus-anchor difference doubles this to $16\eta_0$.
\end{proof}

\subsection{Fallback and downstream transfer}

\begin{proposition}[Exact fallback decomposition]
\label{prop:fallback}
Let $Z$ denote the event that \PSC{} accepts the candidate, and let $L$ be any adverse event defined relative to the anchor such that deploying the anchor itself cannot trigger $L$; for example, $L=\{\Delta_{\AUC}<-0.05\}$. Then
\[
\Pr(L\text{ under \PSC{}})
=
\Pr(Z=1,L\text{ for }S^\star)
=
\Pr(Z=1)\Pr(L\mid Z=1).
\]
\end{proposition}

\begin{proof}
When $Z=0$, \PSC{} deploys $A$ exactly, so every anchor-relative contrast equals zero and $L$ cannot occur. The event $L$ under the deployed policy is therefore identical to the intersection of acceptance and the candidate-level adverse event.
\end{proof}

Proposition~\ref{prop:fallback} explains why low acceptance mechanically reduces realized tail incidence. It does not show that the certificate identifies the right cells. The matched-acceptance diagnostic in the main text addresses that mechanism empirically by comparing \PSC{} with random gates having the same 29/180 acceptance count.

\begin{proposition}[Transfer under a model-discrepancy envelope]
\label{prop:transfer}
Let $h$ be an unseen downstream learner and let $R$ be a loss-like utility metric. Suppose, for the accepted candidate/anchor pair, the downstream contrast obeys
\[
\Delta_R(h;S^\star,A)
\le
\max_{j\le K}\Delta_R(g_j;S^\star,A)+\rho_h.
\]
If all certifier-side bounds in Corollary~\ref{cor:multi} hold with common tolerance $\tau_R$ and radius at most $\varepsilon_R$, then acceptance implies
\[
\Delta_R(h;S^\star,A)
\le \tau_R+\varepsilon_R+\rho_h.
\]
An analogous reversed inequality holds for AUC.
\end{proposition}

\begin{proof}
Substitute the simultaneous certifier-side upper bound into the assumed discrepancy envelope.
\end{proof}

The discrepancy term $\rho_h$ is not known or calibrated in v2.7. Proposition~\ref{prop:transfer} is therefore a statement of what an additional transfer assumption would buy, not a claimed theorem for TabPFN-3.5 or TabICLv2. Their behavior is assessed only through untouched fresh experiments.


\section{Dataset and Split Details}

\begin{table*}[h]
\centering
\caption{Fresh-confirmation tasks after the frozen binary-group filters. All tasks use exact split sizes $3000/500/500/1000/2000$ for $P/C/V/W/T$. Feature counts are before one-hot expansion.}
\label{tab:data}
\small
\begin{tabular}{lrrp{3.25cm}p{3.25cm}}
\toprule
Task & Rows & Features & Target & Protected comparison \\
\midrule
ACS Income & 148252 & 9 & official ACSIncome binary label & White vs Black (RAC1P) \\
ACS Employment & 118051 & 15 & official ACSEmployment binary label & White vs Black (RAC1P) \\
ACS Public Coverage & 64537 & 18 & official ACSPublicCoverage label & White vs Black (RAC1P) \\
ACS Mobility & 45946 & 20 & official ACSMobility binary label & White vs Black (RAC1P) \\
ACS Travel Time & 134774 & 15 & official ACSTravelTime label & White vs Black (RAC1P) \\
ACS Health Insurance & 249131 & 25 & official ACSHealthInsurance label & White vs Black (RAC1P) \\
ACS Income/Poverty & 225874 & 19 & official ratio-task binary label & White vs Black (RAC1P) \\
Law School & 18690 & 10 & first-try bar passage & female vs male \\
Diabetes Hospital & 95309 & 46 & readmission within 30 days & African-American vs Caucasian \\
\bottomrule
\end{tabular}
\end{table*}

Protected attributes are removed from the model features in all tasks.
The ACS tasks use 2023 one-year PUMS records from CO, MI, MN, NJ, and OR and self-contained Folktables-compatible task definitions.

\section{GPU Stratification}

\begin{table}[h]
\centering
\caption{Fresh GPU results by context budget.}
\label{tab:budget}
\small
\begin{tabular}{rrrrrr}
\toprule
$k$ & Accept & $\Delta\eo$ & $\Delta\AUC$ & $\Delta\LL$ & Severe tail \\
\midrule
32 & 13.33\% & -0.000505 & +0.002934 & -0.007055 & 1.111\% \\
64 & 13.33\% & -0.000774 & +0.001608 & -0.002344 & 0.000\% \\
128 & 15.56\% & -0.000363 & -0.000158 & -0.000451 & 0.000\% \\
256 & 22.22\% & -0.000409 & +0.000449 & +0.000211 & 0.000\% \\
\bottomrule
\end{tabular}
\end{table}

The weaker random5 win rates at $k=256$ are one reason we do not claim monotone improvement with context size.

\section{Reproducibility Snapshot}

The frozen method source has SHA256
\begin{center}
\footnotesize\texttt{f02ccffec4bd20225e261bd4a2b88c803ebfbf700e8f53d6b664e42f82dc307d}.
\end{center}
The fresh GPU run used \texttt{tabpfn==9.0.0}, \texttt{tabicl==2.2.0}, PyTorch 2.8.0+cu128, and an NVIDIA RTX 4090 D.
The TabPFN-3.5 checkpoint SHA256 was
\begin{center}
\footnotesize\texttt{ece4d67eadfea42eb0e610df5189bea60cb7f31073d81e9c7a019b76eacf0be3},
\end{center}
and the TabICLv2 checkpoint SHA256 was
\begin{center}
\footnotesize\texttt{bdc7dbd5e4ff21f8f0456fcf90c6b7cdf72dbea960f2d05b19bec19f9b3d4ed0}.
\end{center}
The fresh GPU lattice completed 2880/2880 tasks with zero errors.


\section*{Statements and Declarations}

\paragraph{Funding.}
The author received no specific funding for this work.

\paragraph{Competing interests.}
The authors declare no competing interests.

\paragraph{Author contributions.}
Tao Ran: Conceptualization, Methodology, Software, Validation, Formal analysis, Investigation, Data curation, Visualization, Writing---original draft, Writing---review and editing, and Project administration.

\paragraph{Data availability.}
The study uses public benchmark datasets, including ACS/Folktables, LSAC, and the UCI Diabetes 130-US Hospitals dataset. Dataset provenance and preprocessing are described in the manuscript. Frozen manifests, experiment configurations, context hashes, exact confirmatory result archives, and analysis code are publicly available at \url{https://github.com/Trylaney/psc-bs}.

\paragraph{Code availability.}
The frozen CPU/GPU source, frozen contexts, exact result archives, integrity hashes, audit summaries, and reproduction instructions are publicly available at \url{https://github.com/Trylaney/psc-bs} under the MIT License.

\paragraph{Ethics approval.}
This study performs secondary computational analysis of publicly available benchmark datasets and does not involve new recruitment, intervention, or collection of identifiable human-subject data. Any dataset-specific terms of use and provenance requirements remain applicable.


\begin{thebibliography}{99}

\bibitem[Hu et al.(2024)Hu, Liu, and Du]{hu2024strategic}
Jingyu Hu, Weiru Liu, and Mengnan Du.
\newblock Strategic Demonstration Selection for Improved Fairness in LLM In-Context Learning.
\newblock In \emph{EMNLP}, pages 7460--7475, 2024.

\bibitem[Kenfack et al.(2026)Kenfack, Ebrahimi Kahou, and Aivodji]{kenfack2026fairtabicl}
Patrik Kenfack, Samira Ebrahimi Kahou, and Ulrich Aivodji.
\newblock Towards Fair In-Context Learning with Tabular Foundation Models.
\newblock \emph{Transactions on Machine Learning Research}, 2026.

\bibitem[Chhikara et al.(2025)Chhikara, Ghosh, and Chakraborty]{chhikara2025smite}
Garima Chhikara, Kripabandhu Ghosh, and Abhijnan Chakraborty.
\newblock SMITE: Enhancing Fairness in LLMs through Optimal In-Context Example Selection via Dynamic Validation.
\newblock arXiv:2508.17735, 2025.

\bibitem[Zhang et al.(2025a)Zhang, Xiao, Xiao, Gao, and Zhao]{zhang2025dva}
Qi Zhang, Zhiqing Xiao, Ruixuan Xiao, Lirong Gao, and Junbo Zhao.
\newblock D.Va: Validate Your Demonstration First Before You Use It.
\newblock In \emph{ACL}, pages 2580--2594, 2025.

\bibitem[Gupta et al.(2023)Gupta, Gardner, and Singh]{gupta2023coverage}
Shivanshu Gupta, Matt Gardner, and Sameer Singh.
\newblock Coverage-based Example Selection for In-Context Learning.
\newblock In \emph{Findings of EMNLP}, pages 13924--13950, 2023.

\bibitem[Zhang et al.(2025b)Zhang, Zhang, Li, Wang, Dy, and Zhang]{zzhang2025linear}
Ziniu Zhang, Zhenshuo Zhang, Dongyue Li, Lu Wang, Jennifer Dy, and Hongyang R. Zhang.
\newblock Linear-Time Demonstration Selection for In-Context Learning via Gradient Estimation.
\newblock In \emph{EMNLP}, pages 16459--16477, 2025.

\bibitem[Zhang et al.(2025c)Zhang, Li, Bai, Li, Wang, Lin, and Rong]{jzhang2025gradient}
Jianfei Zhang, Bei Li, Jun Bai, Rumei Li, Yanmeng Wang, Chenghua Lin, and Wenge Rong.
\newblock Selecting Demonstrations for Many-Shot In-Context Learning via Gradient Matching.
\newblock In \emph{Findings of ACL}, pages 11686--11704, 2025.

\bibitem[Hollmann et al.(2025)Hollmann, M{\"u}ller, Purucker, Krishnakumar, K{\"o}rfer, Hoo, Schirrmeister, and Hutter]{hollmann2025tabpfn}
Noah Hollmann, Samuel M{\"u}ller, Lennart Purucker, Arjun Krishnakumar, Max K{\"o}rfer, Shi Bin Hoo, Robin Tibor Schirrmeister, and Frank Hutter.
\newblock Accurate predictions on small data with a tabular foundation model.
\newblock \emph{Nature}, 637:319--326, 2025.

\bibitem[J{\"a}ger et al.(2026)J{\"a}ger et al.]{jaeger2026tabpfn35}
Benjamin J{\"a}ger et al.
\newblock TabPFN-3.5: Technical Report.
\newblock arXiv:2609.17895, 2026.

\bibitem[Qu et al.(2025)Qu, Holzm{\"u}ller, Varoquaux, and Le Morvan]{qu2025tabicl}
Jingang Qu, David Holzm{\"u}ller, Ga{\"e}l Varoquaux, and Marine Le Morvan.
\newblock TabICL: A Tabular Foundation Model for In-Context Learning on Large Data.
\newblock In \emph{ICML}, PMLR 267:50817--50847, 2025.

\bibitem[Qu et al.(2026)Qu, Holzm{\"u}ller, Varoquaux, and Le Morvan]{qu2026tabiclv2}
Jingang Qu, David Holzm{\"u}ller, Ga{\"e}l Varoquaux, and Marine Le Morvan.
\newblock TabICLv2: A better, faster, scalable, and open tabular foundation model.
\newblock arXiv:2602.11139, 2026.

\bibitem[Ding et al.(2021)Ding, Hardt, Miller, and Schmidt]{ding2021retiring}
Frances Ding, Moritz Hardt, John Miller, and Ludwig Schmidt.
\newblock Retiring Adult: New Datasets for Fair Machine Learning.
\newblock In \emph{NeurIPS}, volume 34, 2021.

\bibitem[Wightman(1998)]{wightman1998lsac}
Linda F. Wightman.
\newblock LSAC National Longitudinal Bar Passage Study.
\newblock LSAC Research Report Series, Law School Admission Council, 1998.

\bibitem[Clore et al.(2014)Clore, Cios, DeShazo, and Strack]{clore2014diabetes}
John Clore, Krzysztof Cios, Jonathan DeShazo, and Beata Strack.
\newblock Diabetes 130-US Hospitals for Years 1999--2008.
\newblock UCI Machine Learning Repository, 2014. DOI: 10.24432/C5230J.

\bibitem[Strack et al.(2014)Strack et al.]{strack2014hba1c}
Beata Strack et al.
\newblock Impact of HbA1c Measurement on Hospital Readmission Rates: Analysis of 70,000 Clinical Database Patient Records.
\newblock \emph{BioMed Research International}, 2014:781670, 2014.


\bibitem[Robertson et al.(2025)Robertson, Hollmann, M{\"u}ller, Awad, and Hutter]{robertson2025fairpfn}
Jake Robertson, Noah Hollmann, Samuel M{\"u}ller, Noor Awad, and Frank Hutter.
\newblock FairPFN: A Tabular Foundation Model for Causal Fairness.
\newblock In \emph{ICML}, PMLR 267:51787--51808, 2025.

\bibitem[Kuchibhotla et al.(2022)Kuchibhotla, Kolassa, and Kuffner]{kuchibhotla2022postselection}
Arun K. Kuchibhotla, John E. Kolassa, and Todd A. Kuffner.
\newblock Post-Selection Inference.
\newblock \emph{Annual Review of Statistics and Its Application}, 9:505--527, 2022.

\bibitem[Dwork et al.(2015)Dwork et al.]{dwork2015reusable}
Cynthia Dwork, Vitaly Feldman, Moritz Hardt, Toniann Pitassi, Omer Reingold, and Aaron Roth.
\newblock The reusable holdout: Preserving validity in adaptive data analysis.
\newblock \emph{Science}, 349(6248):636--638, 2015.

\bibitem[Wynn et al.(2026)Wynn, Jazbec, Peris, Khaziev, Liu, Khashabi, and Nalisnick]{wynn2026corrupted}
Andrea Wynn, Metod Jazbec, Charith Peris, Rinat Khaziev, Anqi Liu, Daniel Khashabi, and Eric Nalisnick.
\newblock Controlling the Risk of Corrupted Contexts for Language Models via Early-Exiting.
\newblock In \emph{ICML}, 2026. arXiv:2510.02480.

\bibitem[Angelopoulos et al.(2025)Angelopoulos, Bates, Cand\`es, Jordan, and Lei]{angelopoulos2025ltt}
Anastasios N. Angelopoulos, Stephen Bates, Emmanuel J. Cand\`es, Michael I. Jordan, and Lihua Lei.
\newblock Learn then Test: Calibrating Predictive Algorithms to Achieve Risk Control.
\newblock \emph{The Annals of Applied Statistics}, 2025.

\bibitem[Geifman and El-Yaniv(2017)]{geifman2017selective}
Yonatan Geifman and Ran El-Yaniv.
\newblock Selective Classification for Deep Neural Networks.
\newblock In \emph{NeurIPS}, volume 30, 2017.

\bibitem[Hardt et al.(2016)Hardt, Price, and Srebro]{hardt2016equality}
Moritz Hardt, Eric Price, and Nati Srebro.
\newblock Equality of Opportunity in Supervised Learning.
\newblock In \emph{NeurIPS}, volume 29, 2016.

\bibitem[Hoeffding(1963)]{hoeffding1963probability}
Wassily Hoeffding.
\newblock Probability Inequalities for Sums of Bounded Random Variables.
\newblock \emph{Journal of the American Statistical Association}, 58(301):13--30, 1963.

\bibitem[Cl{\'e}men{\c c}on et al.(2008)Cl{\'e}men{\c c}on, Lugosi, and Vayatis]{clemencon2008ranking}
St{\'e}phan Cl{\'e}men{\c c}on, G{\'a}bor Lugosi, and Nicolas Vayatis.
\newblock Ranking and Empirical Minimization of U-Statistics.
\newblock \emph{The Annals of Statistics}, 36(2):844--874, 2008.

\end{thebibliography}

\end{document}